\documentclass[12pt,a4paper]{article}
\usepackage[utf8]{vietnam}
\usepackage{amsmath,amssymb}
\usepackage{tabularx}
\usepackage{tikz}
\usetikzlibrary{calc,angles,quotes,arrows.meta,patterns}
\usepackage{xcolor}
\usepackage{enumitem}
\usepackage{fancyhdr}
\usepackage[top=1.2cm, bottom=1.2cm, left=1.5cm, right=1.5cm]{geometry}
\usepackage[dethi]{ex_test}
\SAOPTN{widthSA=2.6,heightSA=0.62}

\setlength{\headheight}{15pt}
\setlength{\parskip}{0pt}

% Dòng kẻ chấm mờ chuẩn viết tay
\newcommand{\dongke}[1]{%
  \begingroup%
  \setlength{\parskip}{0pt}%
  \baselineskip=18.5pt%
  \vspace{1.0mm}%
  \foreach \i in {1,...,#1}{%
    \noindent{\color{gray!60}\dotfill}\par%
  }%
  \vspace{0.6mm}%
  \endgroup%
}

\pagestyle{fancy}
\fancyhf{}
\rhead{\textit{Môn Toán 6 -- Đề số 5}}
\lhead{\textit{Bộ đề ôn tập Giữa học kì I}}
\cfoot{\thepage}

\begin{document}

% --- KHUNG TIÊU ĐỀ ĐỀ THI ---
\noindent
\begin{tabularx}{\linewidth}{X r}
\textbf{TRƯỜNG THCS: \dotfill} & \textbf{ĐỀ KIỂM TRA ĐỊNH KỲ GIỮA HỌC KỲ I} \\
\textbf{Họ và tên: \dotfill} & \textbf{NĂM HỌC 2026 -- 2027} \\
\textbf{Lớp: \dotfill \quad SBD: \dotfill} & \textbf{Môn: TOÁN 6} \\
\hline
\end{tabularx}
\vspace{2mm}

\begin{center}
  \textbf{\large ĐỀ ÔN TẬP GIỮA HỌC KỲ I -- ĐỀ SỐ 5}\\
  \textit{Thời gian làm bài: 90 phút (Không kể thời gian phát đề)}
\end{center}
\vspace{1mm}

% --- PHẦN I ---
\noindent\textbf{\large PHẦN I. CÂU TRẮC NGHIỆM NHIỀU PHƯƠNG ÁN LỰA CHỌN (3,0 điểm)}

\textit{Thí sinh trả lời từ câu 1 đến câu 12. Mỗi câu hỏi thí sinh chỉ chọn một phương án duy nhất.}

\vspace{1mm}
\begin{ex}
Cho $A$ là tập hợp các số tự nhiên nhỏ hơn $5$. Tập hợp $A$ được viết như thế nào?
\choice
{$A = \{1; 2; 3; 4\}$}
{\True $A = \{0; 1; 2; 3; 4\}$}
{$A = \{1; 2; 3; 4; 5\}$}
{$A = \{0; 1; 2; 3; 4; 5\}$}
\end{ex}

\begin{ex}
Kết quả của phép tính $5^2 \cdot 5$ là:
\choice
{$5$}
{$5^2$}
{\True $5^3$}
{$25^2$}
\end{ex}

\begin{ex}
Lũy thừa $2^4$ có giá trị bằng lũy thừa nào sau đây?
\choice
{\True $4^2$}
{$3^3$}
{$5^1$}
{$2^2 + 2^2$}
\end{ex}

\begin{ex}
Số nào sau đây chia hết cho $5$ nhưng không chia hết cho $3$?
\choice
{\True $325$}
{$240$}
{$225$}
{$600$}
\end{ex}

\begin{ex}
Số nào dưới đây là một bội chung của $3, 4, 5$?
\choice
{\True $120$}
{$190$}
{$145$}
{$150$}
\end{ex}

\begin{ex}
Tam giác $ABC$ là tam giác đều, khi đó:
\choice
{$AB < BC < CA$}
{$AB > BC > CA$}
{Độ dài ba cạnh $AB, BC, CA$ đôi một khác nhau}
{\True $AB = BC = CA$}
\end{ex}

\begin{ex}
Chọn phát biểu đúng trong các phát biểu sau:
\choice
{$0$ là bội của mọi số tự nhiên}
{Các số nguyên tố nhỏ hơn 12 là $1; 2; 5; 7$}
{$1$ là số nguyên tố nhỏ nhất}
{\True Số $2$ là số nguyên tố chẵn duy nhất}
\end{ex}

\begin{ex}
Công thức nào sau đây thể hiện tính chất phân phối của phép nhân đối với phép cộng?
\choice
{\True $a(b + c) = ab + ac$}
{$(a + b) + c = a + (b + c)$}
{$ab = ba$}
{$(ab)c = a(bc)$}
\end{ex}

\begin{ex}
Phát biểu nào dưới đây là sai?
\choice
{Hình vuông có hai đường chéo vuông góc với nhau}
{Hình lục giác đều có sáu góc bằng nhau}
{\True Hình tam giác đều có ba đường chéo bằng nhau}
{Hình vuông có bốn góc bằng nhau và bằng $90^\circ$}
\end{ex}

\begin{ex}
Cho hai tập hợp $M = \{a; 5; b; c\}$ và $N = \{d; 3; e; 7\}$. Khẳng định nào sau đây là sai?
\choice
{\True $5 \in N$}
{$3 \in N$}
{$a \in M$}
{$c \in M$}
\end{ex}

\begin{ex}
Trong các số sau, số nào chia hết cho $5$?
\choice
{$553$}
{\True $12340$}
{$578$}
{$786$}
\end{ex}

\begin{ex}
Hình nào trong các hình sau có hai cạnh đối song song và hai đường chéo bằng nhau?
\choice
{Hình bình hành}
{\True Hình thang cân}
{Hình thoi}
{Hình tam giác đều}
\end{ex}


\vspace{1mm}
\noindent\textbf{BẢNG TRẢ LỜI TRẮC NGHIỆM PHẦN I:}
\begin{center}
\begin{tabular}{|c|c|c|c|c|c|c|c|c|c|c|c|c|}
\hline
\textbf{Câu} & \textbf{1} & \textbf{2} & \textbf{3} & \textbf{4} & \textbf{5} & \textbf{6} & \textbf{7} & \textbf{8} & \textbf{9} & \textbf{10} & \textbf{11} & \textbf{12} \\
\hline
\textbf{Chọn} & \rule{0pt}{0.65cm} & & & & & & & & & & & \\
\hline
\end{tabular}
\end{center}
\vspace{2mm}

\setcounter{ex}{0}
\noindent\textbf{\large PHẦN II. CÂU TRẮC NGHIỆM ĐÚNG SAI (2,0 điểm)}

\textit{Thí sinh trả lời câu 1 và câu 2. Trong mỗi ý a), b), c), d) ở mỗi câu, thí sinh chọn đúng hoặc sai.}

\vspace{1mm}
\begin{ex}
Xét tính đúng hay sai của các mệnh đề sau:
\choiceTF
{\True $11$ là số nguyên tố}
{$15$ là số nguyên tố}
{\True $17$ là số nguyên tố}
{$21$ là số nguyên tố}
\end{ex}

\begin{ex}
Cho $15$ số tự nhiên đầu tiên (từ $0$ đến $14$). Xét tính đúng hay sai của các mệnh đề sau:
\choiceTF
{\True Trong các số đã cho, có $2$ số chia hết cho cả $2$ và $5$}
{Trong các số đã cho, có $5$ số chia hết cho $9$}
{\True Trong các số đã cho, có $6$ số là số nguyên tố}
{\True Trong các số đã cho, có $5$ số là bội của $3$}
\end{ex}


\noindent\textbf{BẢNG TRẢ LỜI PHẦN II:}
\begin{center}
\begin{tabular}{|c|c|c|c|c||c|c|c|c|c|}
\hline
\textbf{Câu 1} & \textbf{a)} & \textbf{b)} & \textbf{c)} & \textbf{d)} & \textbf{Câu 2} & \textbf{a)} & \textbf{b)} & \textbf{c)} & \textbf{d)} \\
\hline
\textbf{Đ/S} & \rule{0pt}{0.6cm} & & & & \textbf{Đ/S} & & & & \\
\hline
\end{tabular}
\end{center}
\vspace{2mm}

\setcounter{ex}{0}
\noindent\textbf{\large PHẦN III. CÂU HỎI TRẢ LỜI NGẮN (2,0 điểm)}

\textit{Thí sinh trả lời câu 1 và câu 2. Điền kết quả số vào ô trống tương ứng.}

\vspace{1mm}
\begin{ex}
Tính diện tích của một hình vuông có chu vi bằng $36\text{ cm}$ (đơn vị: $\text{cm}^2$). \shortans{$81$}
\end{ex}

\begin{ex}
Tìm số tự nhiên $x$ nhỏ nhất khác $0$ sao cho $x \ \vdots \ 12$ và $x \ \vdots \ 18$. \shortans{$36$}
\end{ex}


\setcounter{ex}{0}
\noindent\textbf{\large PHẦN IV. TỰ LUẬN (3,0 điểm)}

\textit{Thí sinh trình bày chi tiết lời giải vào giấy làm bài.}

\vspace{1mm}
\begin{ex} (1.0 điểm)
Thực hiện phép tính:
\begin{enumerate}[label=\textbf{\alph*)}, leftmargin=0.8cm, itemsep=2pt]
\item $3 \cdot 5^2 - 16 : 2^2$
\item $20 - [30 - (5 - 1)^2]$
\end{enumerate}
\end{ex}
\vspace{1.5mm}

\begin{ex} (0.75 điểm)
Tìm số tự nhiên $x$, biết:
\begin{enumerate}[label=\textbf{\alph*)}, leftmargin=0.8cm, itemsep=2pt]
\item $x - 15 = 2025$
\item $39 + (x - 47) = 160$
\end{enumerate}
\end{ex}
\vspace{1.5mm}

\begin{ex} (0.75 điểm)
Bạn Ánh đi tập gym mỗi ngày. Lượng calo tiêu hao trung bình khi tập mỗi buổi gồm:\
$\bullet$ Chạy bộ $30$ phút tiêu hao $240$ calo.\
$\bullet$ Nhảy dây $15$ phút tiêu hao $150$ calo.\
$\bullet$ Gập bụng $10$ phút tiêu hao $60$ calo.
\begin{enumerate}[label=\textbf{\alph*)}, leftmargin=0.8cm, itemsep=2pt]
\item Tính tổng lượng calo tiêu hao trong một buổi tập của bạn Ánh.
\item Biết người bình thường cần nạp khoảng $2200$ calo mỗi ngày để giữ cân và cần tiêu hao $7700$ calo trong một tuần để giảm $1\text{ kg}$ mỡ. Nếu mỗi ngày bạn Ánh chỉ tập các bài tập như trên mà vẫn ăn uống đủ $2200$ calo thì sau $1$ tuần bạn ấy đã đủ để giảm $1\text{ kg}$ mỡ chưa?
\end{enumerate}
\end{ex}
\vspace{1.5mm}

\begin{ex} (0.5 điểm)
\begin{enumerate}[label=\textbf{\alph*)}, leftmargin=0.8cm, itemsep=2pt]
\item Tìm số tự nhiên $n$ sao cho $(5n + 14) \ \vdots \ (n + 2)$.
\end{enumerate}
\end{ex}
\vspace{1.5mm}


\begin{center}
\rule{4cm}{0.4pt}\\
\textit{--- HẾT ---}
\end{center}
\end{document}
